\section{Actual monomial support on the same boundary-incidence stratum}

\begin{lemma}[The whole monomial part is a unit precisely when every incident exponent vanishes]
\label{hir:sharp-actual-monomial-unit-incidence}
\uses{hir:monomial-part-along,hir:exponent-along}
\lean{TranscendentalNakamaye.Pending.HironakaSharpActualTraceIncidence.one_mem_monomialPartAlong_iff}
\leanok
At a point of the actual carrier, the constant germ one belongs to the
whole monomial part if and only if every boundary member through that
point has exponent zero. This includes the empty boundary and zero
exponents; no chosen local equation or unit certificate is supplied.
\end{lemma}
\begin{proof}
For an incident member with positive exponent, every germ of the
corresponding positive power of its vanishing ideal vanishes at the
point: expand its outer product span and evaluate, using the vanishing
of the last factor. Every relative finite span of such germs also
vanishes there, since the point belongs to the carrier. Thus membership
of the constant one in the monomial part would give one equal to zero.
Conversely, if every incident exponent is zero, each required ideal
power is the unit ideal. The constant one is its own single relative
generator, which satisfies every clause of the actual monomial-part
definition. No condition is imposed by a nonincident member.
\end{proof}

\begin{theorem}[Actual monomial-unit status is locally constant on a fixed incidence stratum]
\label{hir:sharp-actual-monomial-unit-stable-boundary-stratum}
\uses{hir:sharp-actual-monomial-unit-incidence,hir:exponent-along}
\lean{TranscendentalNakamaye.Pending.HironakaSharpActualTraceIncidence.eventually_monomial_unit_of_stable_incidence}
\leanok
In positive ambient dimension, let the actual smooth carrier be transverse to an original SNC boundary,
and let the actual ideal be coherent and nonzero along the carrier.
On a subset of the carrier whose whole boundary-incidence profile is
locally constant at a point, membership of one in the whole monomial
part is locally constant at that same point.
\end{theorem}
\begin{proof}
For each of the finitely many members incident at the point, the
genuine divisor-exponent theorem gives constancy of its actual exponent
on the nearby trace in the carrier. Intersect these finitely many real
neighbourhood conditions with the given whole incidence neighbourhood.
A nearby point of the subset has exactly the same incident members,
and the exponent on each of them is the same. The preceding unit
criterion therefore has exactly the same truth value at both points.
Members absent at the original point remain absent and require no
exponent comparison.
\end{proof}

\begin{lemma}[Restriction to the actual carrier preserves the whole zero locus there]
\label{hir:sharp-actual-restricted-whole-zero-locus}
\uses{ana:restricted-calculus}
\lean{TranscendentalNakamaye.Pending.HironakaSharpActualTraceIncidence.zeroLocus_restrictOn_iff}
\leanok
At a point of the carrier, the whole restricted extension of a coherent
ideal vanishes if and only if the original whole ideal vanishes.
\end{lemma}
\begin{proof}
Every original germ belongs to the restricted extension, so vanishing of
the extension implies vanishing of the original ideal. Conversely every
germ of the extension agrees on the carrier, near the point, with a finite
holomorphic combination of original ideal germs. At a point of the carrier
where all original germs vanish, evaluate this actual finite equality.
Every summand vanishes, hence so does the extension germ. This proves both
directions for the entire ideal.
\end{proof}

\begin{theorem}[A genuine whole trace factor and the two actual profiles force constant old-member incidence]
\label{hir:sharp-actual-factor-forces-trace-incidence}
\uses{hir:sharp-actual-monomial-unit-stable-boundary-stratum,hir:sharp-actual-restricted-whole-zero-locus,hir:sharp-actual-full-maximum-original-member-orders}
\lean{TranscendentalNakamaye.Pending.HironakaSharpActualTraceIncidence.eventually_trace_incidence_of_factor}
\leanok
In positive ambient dimension, let a smooth carrier be transverse to an actual SNC born boundary, and let
$H$ be an actual closed smooth old member whose trace ideal is nonzero
along the carrier. Suppose the current whole trace has the actual factor
$I_H|_N=(D J)|_N$, where $D$ is its actual monomial part and $J$ is coherent.
This equality is the explicitly retained current-stage induction hypothesis.
On a subset of the carrier, assume its whole born-boundary incidence
profile and the unit profile of this same $J$ are locally constant at a
point. Then incidence with the actual old member $H$ is locally constant
on that same subset near that point. No centre is replaced by pieces.
\end{theorem}
\begin{proof}
The true monomial-part theorem makes $D$ coherent. The preceding actual
exponent and incidence argument makes membership of one in $D$ locally
constant. For any coherent ideal, failure of unit membership is equivalent
to vanishing of the entire ideal: a germ with nonzero value has a genuine
holomorphic inverse near the point, which puts one in the ideal; and one
cannot vanish there. Thus the whole zero-locus profiles of $D$ and $J$ are
both locally constant on the same subset near the chosen point.

The zero locus of the actual product $D J$ is their union. Indeed every
product germ vanishes if either factor ideal vanishes, whereas nonvanishing
germs in both factors give a nonvanishing product. The preceding restriction
lemma preserves this union at every point of the carrier. Substitute the
explicit whole trace factor and the actual closed vanishing-ideal dictionary
to identify incidence with $H$ with that union. Intersect the two actual
profile neighbourhoods. The truth value of their union, hence of incidence
with $H$, is constant there on the original unchanged subset.
\end{proof}
